\documentclass[../m00-session-01.tex]{subfiles}
\begin{document}
\TopicStandaloneTitle{02}{Make the workflow explicit}{Foundations: Python, Cloud and a reproducible workspace}
\TopicDivider{02}{Make the workflow explicit}{Use an environment, version control and clear responsibility boundaries.}

\begin{frame}[fragile]{First checkpoint: prove the environment is explicit}
  \begin{CodeBlock}[numbers=left]{Bash}
python --version
python -m venv .venv
source .venv/bin/activate
python -m pip install pandas scikit-learn
git init
git status
  \end{CodeBlock}
  \begin{WarningBox}{Do not treat an activated environment as documentation}
    Record dependencies in a requirements file or equivalent lockfile so the session can be reproduced later.
  \end{WarningBox}
\end{frame}

\begin{frame}{A practical responsibility map}
  \centering
  \begin{tikzpicture}[node distance=1.25cm]
    \node[upgrad/service] (git) {Git\\code history};
    \node[upgrad/data,right=of git] (s3) {S3\\data + artefacts};
    \node[upgrad/model,right=of s3] (python) {Python\\analysis};
    \draw[upgrad/flow] (git) -- node[above,font=\scriptsize,text=UpGradSlate]{clone} (python);
    \draw[upgrad/flow] (s3) -- node[above,font=\scriptsize,text=UpGradSlate]{read / write} (python);
  \end{tikzpicture}
  \vspace{.35cm}
  \begin{TakeawayBox}{Reproducibility begins before modelling}
    A clear project boundary makes it possible to tell which code, data and credentials created an analysis.
  \end{TakeawayBox}
\end{frame}
\end{document}
